% \iffalse meta-comment
%<*internal>
\def\nameofplainTeX{plain}
\ifx\fmtname\nameofplainTeX\else
  \expandafter\begingroup
\fi
%</internal>
%<*install>
\input docstrip.tex
\keepsilent
\askforoverwritefalse
\preamble
----------------------------------------------------------------
exam-gradescope --- Gradescope outlines, rubrics and answer keys
Released under the LaTeX Project Public License v1.3c or later
See http://www.latex-project.org/lppl.txt
----------------------------------------------------------------

\endpreamble
\postamble

Copyright (C) 2026 by Matthew Leingang

This work may be distributed and/or modified under the
conditions of the LaTeX Project Public License (LPPL), either
version 1.3c of this license or (at your option) any later
version.  The latest version of this license is in the file:

http://www.latex-project.org/lppl.txt

This work is "maintained" (as per LPPL maintenance status) by
Matthew Leingang.

This work consists of the file  exam-gradescope.dtx
and the derived files           exam-gradescope.ins,
                                exam-gradescope.pdf,
                                exam-gradescope.sty and
                                exam-gradescope.lua.

\endpostamble
\usedir{tex/lualatex/exam-gradescope}
\generate{
  \file{\jobname.sty}{\from{\jobname.dtx}{package}}
}
\def\MetaPrefix{--}
\preamble
----------------------------------------------------------------
exam-gradescope --- Gradescope outlines, rubrics and answer keys
Released under the LaTeX Project Public License v1.3c or later
See http://www.latex-project.org/lppl.txt
----------------------------------------------------------------

\endpreamble
\postamble

Copyright (C) 2026 by Matthew Leingang

This work may be distributed and/or modified under the
conditions of the LaTeX Project Public License (LPPL), either
version 1.3c of this license or (at your option) any later
version.  The latest version of this license is in the file:

http://www.latex-project.org/lppl.txt

This work is "maintained" (as per LPPL maintenance status) by
Matthew Leingang.

This work consists of the file  exam-gradescope.dtx
and the derived files           exam-gradescope.ins,
                                exam-gradescope.pdf,
                                exam-gradescope.sty and
                                exam-gradescope.lua.

\endpostamble
\generate{
  \file{\jobname.lua}{\from{\jobname.dtx}{lua}}
}
%</install>
%<install>\endbatchfile
%<*internal>
\usedir{source/latex/exam-gradescope}
\generate{
  \file{\jobname.ins}{\from{\jobname.dtx}{install}}
}
\nopreamble\nopostamble
\usedir{doc/latex/exam-gradescope}
\ifx\fmtname\nameofplainTeX
  \expandafter\endbatchfile
\else
  \expandafter\endgroup
\fi
%</internal>
%<*package>
\NeedsTeXFormat{LaTeX2e}[2023/06/01]
\ProvidesPackage{exam-gradescope}
  [2026/10/30 v0.1 Gradescope outlines, rubrics and answer keys for the exam class]
%</package>
%<*driver>
\ProvidesFile{exam-gradescope.dtx}
  [2026/10/30 v0.1 Gradescope outlines, rubrics and answer keys for the exam class]
\documentclass{ltxdoc}
\usepackage{hologo}
\usepackage[hidelinks]{hyperref}
\providecommand\LuaLaTeX{\hologo{LuaLaTeX}}
\newcommand\option[1]{\texttt{#1}}
\newcommand\env[1]{\textsf{#1}}
\EnableCrossrefs
\CodelineIndex
\RecordChanges
\begin{document}
  \DocInput{\jobname.dtx}
  \PrintChanges
  \PrintIndex
\end{document}
%</driver>
% \fi
%
%\GetFileInfo{\jobname.dtx}
%
%\title{^^A
%  \textsf{exam-gradescope} --- Gradescope outlines, rubrics and answer keys for the exam class\thanks{^^A
%    This file describes version \fileversion, last revised \filedate.^^A
%  }^^A
%}
%\author{^^A
%  Matthew Leingang^^A
%}
%\date{Released \filedate}
%
%\maketitle
%
%\changes{v0.1}{2026/10/30}{First version}
%
% \begin{abstract}
%   The \textsf{exam-gradescope} package records the questions of a document
%   made with the \textsf{exam} class while it is typeset with \LuaLaTeX, and
%   writes the files that Gradescope (\url{https://www.gradescope.com}) needs
%   to grade it: an outline of where each question is on the page, a rubric
%   for each question, an answer key for a bubble sheet, and a manifest of the
%   questions for item analysis.
% \end{abstract}
%
% \tableofcontents
%
% \section{Introduction}
%
% Grading a paper exam on Gradescope starts with an \emph{outline}: a box,
% drawn by hand on the blank exam, around each question and around the
% student's name and ID. Each question then needs a \emph{rubric}. For a
% multiple-choice question that means copying the right answer, which differs
% from version to version, into the rubric by hand. Questions answered on a
% bubble sheet need an answer key instead.
%
% Everything this takes is already in the \LaTeX{} source: where each question
% falls on the page, how many points it is worth, which choice is correct.
% This package writes it down while the document is typeset:
%
% \begin{description}
%   \item[\meta{jobname}.outline.json] the question regions and the name and
%     ID regions, as percentages of the page;
%   \item[\meta{jobname}.rubric.json] a rubric for every question: made from
%     the correct choices for multiple-choice and check-all-that-apply
%     questions, and written in the source (see section~\ref{sec:rubrics}) for
%     free-response ones;
%   \item[\meta{jobname}.key.csv] the answer key of the questions answered on
%     a bubble sheet;
%   \item[\meta{jobname}.json] a manifest of all the questions: number, a
%     name that is the same on every version, points, type and answer.
% \end{description}
% All of them are made from the same run as the PDF. In particular the answer
% key and rubrics come from the blank exam, after any shuffling of choices,
% not from a separate solutions run.
%
% Gradescope has no public interface for uploading outlines and rubrics. The
% command \texttt{edubag gradescope client setup-exam} of the Python package
% \textsf{edubag} (\url{https://github.com/nyumathclinic/python-edubag}) reads
% these files and sets up a multi-version exam on Gradescope; the
% \texttt{gradescope-tampermonkey} rubric importer
% (\url{https://github.com/nyumathclinic/gradescope-tampermonkey}) reads the
% rubric files in a browser.
%
% \section{Requirements}
%
% \begin{itemize}
%   \item \LuaLaTeX{} (the package records positions with Lua at shipout);
%   \item the \textsf{exam} class, or a class built on it;
%   \item a \LaTeX{} kernel from 2023 or later.
% \end{itemize}
% The package works with \textsf{exam-randomizechoices}, loaded before or after
% it.
%
% \section{Usage}
%
% Load the package after the \textsf{exam} class:
% \begin{verbatim}
% \documentclass{exam}
% \usepackage{exam-gradescope}
% \end{verbatim}
% Nothing else is needed for multiple-choice questions and for the outline of
% the questions. The files are written at the end of the run, next to the PDF.
%
% \subsection{Where questions go}
%
% Each question is \emph{multiple choice} if it has a \env{choices} or
% \env{oneparchoices} environment with a |\CorrectChoice|,
% \emph{check all that apply} if it has a \env{checkboxes} or
% \env{oneparcheckboxes} environment, and \emph{free response} otherwise.
%
% Free-response questions are graded on Gradescope. Multiple-choice and
% check-all-that-apply questions go where the option \option{mc} says:
% \begin{description}
%   \item[\texttt{mc=gradescope}] (the default) they are graded on Gradescope
%     like the others, with rubrics made from their correct choices;
%   \item[\texttt{mc=bubble}] they are answered on a bubble sheet: they go in
%     the answer key, and are left out of the outline and rubrics.
% \end{description}
% So a quiz graded entirely on Gradescope needs no option, and an exam whose
% multiple-choice part is on a bubble sheet uses
% \begin{verbatim}
% \usepackage[mc=bubble]{exam-gradescope}
% \end{verbatim}
% Its outline and rubrics then have only the free-response questions. These are
% numbered 1, 2, \dots{} on Gradescope, so each is titled after its printed
% number (``Question 29''). The outline's page numbers are those of the whole
% PDF, which is the template to upload.
%
% \DescribeMacro{\gradescopequestion}
% |\gradescopequestion{|\meta{keys}|}|, anywhere in a question, changes how
% that question is treated:
% \begin{description}
%   \item[\texttt{type=}\meta{mc, cata or frq}] its type;
%   \item[\texttt{destination=}\meta{bubble, gradescope or none}] where it is
%     graded (\texttt{none} leaves it out of everything but the manifest);
%   \item[\texttt{title=}\meta{text}] its title on Gradescope;
%   \item[\texttt{id=}\meta{name}] its name in the manifest;
%   \item[\texttt{points=}\meta{number}] its points, when the document does
%     not give them to \textsf{exam}.
% \end{description}
% In a part, after |\part|, the keys \texttt{title} and \texttt{points} apply
% to the part. A question whose parts have no points of their own has its
% points split evenly among them, with a warning.
%
% \subsection{Names of questions}
%
% The manifest gives each question a name, so that the same question can be
% found on every version even when the questions are shuffled. The name is
% the first |\label| in the question:
% \begin{verbatim}
% \question[2]\label{binomial-sum} ...
% \end{verbatim}
% Labels with an |@| are ignored: \textsf{exam} and
% \textsf{exam-randomizechoices} make those for their own use. A question with
% no label is named by its number.
%
% \subsection{The name and ID regions}
%
% \DescribeMacro{\GradescopeRegion}
% |\GradescopeRegion{|\meta{name or sid}|}{|\meta{material}|}| typesets
% \meta{material} in a box and records where it is, as the region of the
% student's name (\texttt{name}) or ID (\texttt{sid}). Wrap the boxes on the
% cover page in it:
% \begin{verbatim}
% \GradescopeRegion{name}{\fbox{\parbox{3in}{Name:\vspace{2em}}}}\hfill
% \GradescopeRegion{sid}{\fbox{\parbox{2in}{ID:\vspace{2em}}}}
% \end{verbatim}
%
% \subsection{Rubrics}\label{sec:rubrics}
%
% A multiple-choice question gets the rubric
% \begin{quote}
%   \begin{tabular}{rl}
%     \meta{points} & Correct \\
%     0 & The correct response is \textbf{B}
%   \end{tabular}
% \end{quote}
% with positive scoring. A check-all-that-apply question with five choices gets
% negative scoring and the items
% \begin{quote}
%   \begin{tabular}{rl}
%     \meta{points}/5 & A \textbf{should} be selected \\
%     \meta{points}/5 & B \textbf{should not} be selected \\
%     \multicolumn{2}{c}{\dots} \\
%     0 & Correct
%   \end{tabular}
% \end{quote}
% Every rubric has a floor and a ceiling, so a score stays between 0 and the
% question's points.
%
% \DescribeEnv{rubric}
% A free-response question's rubric is written in a \env{rubric} environment
% inside the question (or part):
% \begin{verbatim}
% \begin{question}[6]\label{comb-proof}
%   Give a combinatorial proof that ...
%   \begin{rubric}[scoring=positive]
%     \begin{rubricgroup}[select one]{What is counted}
%       \item[2] Clearly identifies the collection being counted
%       \item[1] Identifies a collection, but vaguely
%       \item[0] No collection identified
%     \end{rubricgroup}
%     \item[0] No response, or no creditable work shown
%   \end{rubric}
% \end{question}
% \end{verbatim}
% Each |\item[|\meta{points}|]| is a rubric item; a \env{rubricgroup} is a
% group of items, which with \texttt{select one} are mutually exclusive. The
% options of \env{rubric} are \texttt{scoring=positive} or
% \texttt{scoring=negative}, \texttt{floor=}\meta{true or false} and
% \texttt{ceiling=}\meta{true or false}; the defaults are positive, true and
% true.
%
% The descriptions are read as written, not typeset: they go to Gradescope,
% whose rubrics use Markdown with \TeX{} math. Math in |\(...\)|, |\[...\]| or
% |$...$| becomes |$$...$$|, |\textbf{...}| becomes |**...**|, and
% |\emph{...}| and |\textit{...}| become |*...*|. Other commands are left
% alone, so keep to ones that MathJax knows inside math.
%
% \DescribeMacro{\DeclareRubric}
% \DescribeMacro{\userubric}
% Many questions share a rubric. |\DeclareRubric[|\meta{options}|]{|^^A
% \meta{name}|}{|\meta{items}|}| declares one, in the preamble, with points as
% fractions of the question's points:
% \begin{verbatim}
% \DeclareRubric{basic}{
%   \item[1]   Correct
%   \item[5/6] Mostly correct, but small mistakes
%   \item[2/3] Major mistakes
%   \item[1/3] Nice try
%   \item[0]   No response, or no creditable work shown}
% \end{verbatim}
% and |\userubric{|\meta{name}|}| in a question uses it, scaled to the
% question's points. An optional argument changes it for that question:
% \begin{verbatim}
% \question[3] ...
% \userubric{basic}[
%   \reword{Correct}{Correct, and explained}
%   \add[before=Major mistakes]{2}{Correct, but not explained}
%   \remove{Nice try}]
% \end{verbatim}
% |\reword{|\meta{old}|}{|\meta{new}|}| rewords the item (or group) whose
% description starts with \meta{old}; |\add[before=|\meta{item}|]| or
% |\add[after=|\meta{item}|]{|\meta{points}|}{|\meta{text}|}| adds an item
% (at the end without an option), worth \meta{points} points, not a fraction;
% |\remove{|\meta{item}|}| removes one.
%
% With the option \texttt{default rubric=}\meta{name}, every free-response
% question (or part) without a rubric of its own gets that one.
%
% \subsection{Options}
%
% Options are given when loading the package or, except for \option{debugregions},
% \DescribeMacro{\gradescopesetup} with |\gradescopesetup{|\meta{keys}|}|.
% \begin{description}
%   \item[\texttt{mc=}\meta{gradescope or bubble}] where multiple-choice and
%     check-all-that-apply questions are graded (default \texttt{gradescope});
%   \item[\texttt{showrubric}] print each rubric under its question when the
%     \textsf{exam} class shows answers (with its \texttt{answers} option), for
%     the graders' copy;
%   \item[\texttt{default rubric=}\meta{name}] the rubric of free-response
%     questions that have none;
%   \item[\texttt{margin=}\meta{length}] how far question regions extend left
%     and right of the text area (default 18pt);
%   \item[\texttt{outline}, \texttt{rubric}, \texttt{key}, \texttt{manifest}]
%     \texttt{=false} to not write that file;
%   \item[\texttt{debugregions}] draw the regions of the previous run on each
%     page, to check them. (Regions are known only once the whole document has
%     been typeset, so this shows the last run's.)
% \end{description}
%
% \section{How regions are found}
%
% A question's region starts where the question starts and runs to where the
% next one starts, across pages if need be, within the text area widened by
% \option{margin}. The last question runs to the foot of its page, where the
% space for its answer is. A part's region runs to the next part, or to the
% end of the question. Slivers shorter than 12pt, left when a question starts
% at the very foot of a page, are dropped.
%
% Each start is recorded with |\latelua|, which gives the exact position on the
% page when the page is shipped out; so the files are right after a single
% run, whatever happens to the page breaks.
%
% \section{The files}
%
% The outline is the body Gradescope's outline editor sends to its server:
% \begin{verbatim}
% {"identification_regions": {
%    "name": {"page_number": 1, "x1": 11.8, "x2": 48.9, "y1": 18.4, "y2": 25.1},
%    "sid":  {"page_number": 1, "x1": 49.2, "x2": 86.4, "y1": 18.4, "y2": 25.1}},
%  "question_data": [
%    {"title": "", "weight": 2,
%     "crop_rect_list": [{"page_number": 1, "x1": 8.8, "x2": 91.2,
%                         "y1": 37.8, "y2": 49.1}]},
%    {"title": "", "children": [{"title": "(a)", "weight": 3, ...}, ...]}]}
% \end{verbatim}
% Coordinates are percentages of the page from its top left corner.
%
% The rubric file is in the format of the \texttt{gradescope-tampermonkey}
% importer: |{"questions": [{"question": "1", "scoring_type": ..., "floor": ...,|
% |"ceiling": ..., "rubric_items": [...], "rubric_item_groups": [...]}]}|,
% where \texttt{question} is the question's number on Gradescope (|"2.1"| for
% the first part of the second question).
%
% The answer key has a header line |Question,Answers,Points,Type| and a line
% such as |3,B,2.0,Exact| for each bubble-sheet question.
%
% The manifest is |{"items": [...]}| with, for each question, its
% \texttt{number}, \texttt{id}, \texttt{points}, \texttt{type} (\texttt{mc},
% \texttt{cata} or \texttt{frq}), \texttt{destination}, \texttt{answer} (the
% correct letters, if any) and \texttt{parts} (a list of |{"points": ...}|).
%
% \section{Limitations}
%
% \begin{itemize}
%   \item Subparts are not regions of their own; their points count towards
%     their part.
%   \item The text area is taken from the page layout at the start of the
%     document. Pages with a different layout (\textsf{exam}'s |\extrawidth|,
%     say) get regions that may be a little narrow.
%   \item Rubric descriptions are converted to Markdown by a few simple rules
%     (above); anything else is passed on as written.
%   \item The bubble-sheet answer key is for multiple-choice questions: a
%     check-all-that-apply question there is written with all its correct
%     letters.
% \end{itemize}
%
% \StopEventually{}
%
% \section{Implementation}
%
% \subsection{The package}
%
%    \begin{macrocode}
%<*package>
\ifdefined\directlua\else
  \PackageError{exam-gradescope}{This package needs LuaLaTeX}{Compile with lualatex.}
  \expandafter\endinput
\fi
\ifdefined\questions\else
  \PackageError{exam-gradescope}{This package needs the exam class}{Use \string\documentclass{exam}, or a class built on it.}
  \expandafter\endinput
\fi

\RequirePackage{etoolbox}
\directlua{examgradescope = require("exam-gradescope")}
%</package>
%    \end{macrocode}
%
% \subsubsection{Options}
%
%    \begin{macrocode}
%<*package>
\newif\ifeg@debug
\newif\ifeg@show
\newif\ifeg@outline \eg@outlinetrue
\newif\ifeg@rubric \eg@rubrictrue
\newif\ifeg@key \eg@keytrue
\newif\ifeg@manifest \eg@manifesttrue
\newcommand\eg@option[2]{\directlua{examgradescope.options["#1"] = #2}}
\DeclareKeys[examgradescope]{
  mc.code = \ifstrequal{#1}{bubble}{\eg@option{mc}{"bubble"}}{%
    \ifstrequal{#1}{gradescope}{\eg@option{mc}{"gradescope"}}{%
      \PackageError{exam-gradescope}{mc should be bubble or gradescope, not '#1'}{}}},
  debugregions.if = eg@debug,
  debugregions.usage = load,
  showrubric.if = eg@show,
  margin.code = \eg@option{margin}{\number\dimexpr#1\relax},
  default~rubric.code = \eg@option{default_rubric}{"\luaescapestring{#1}"},
  outline.if = eg@outline,
  rubric.if = eg@rubric,
  key.if = eg@key,
  manifest.if = eg@manifest,
}
\ProcessKeyOptions[examgradescope]
\NewDocumentCommand\gradescopesetup{m}{\SetKeys[examgradescope]{#1}}
%</package>
%    \end{macrocode}
%
% \subsubsection{Questions and parts}
%
% A position mark, resolved at shipout. |\latelua| expands its argument then, so the arguments are expanded now.
%
%    \begin{macrocode}
%<*package>
\newcommand\eg@mark[2]{%
  \expanded{\noexpand\latelua{examgradescope.mark("#1", #2)}}}

\newcount\eg@index
\newcount\eg@part
%</package>
%    \end{macrocode}
%
% exam defines |\process@question| and |\process@part| inside the questions and parts environments, so they are patched as those begin.
%
%    \begin{macrocode}
%<*package>
\apptocmd\questions{%
  \pretocmd\process@question{%
    \ifhmode\par\fi
    \global\eg@index=\directlua{tex.sprint(examgradescope.start_question(\the\numexpr\value{question}+1\relax))}\relax
    \eg@mark{question}{\the\eg@index}%
  }{}{\PackageError{exam-gradescope}{Could not patch exam's \string\process@question}{}}%
}{}{\PackageError{exam-gradescope}{Could not patch the questions environment}{}}
\AddToHook{env/questions/end}{\ifhmode\par\fi\eg@mark{end}{0}}

\apptocmd\parts{%
  \pretocmd\process@part{%
    \ifhmode\par\fi
    \global\eg@part=\directlua{tex.sprint(examgradescope.start_part(\the\numexpr\value{partno}+1\relax))}\relax
    \eg@mark{part}{\the\eg@index, \the\eg@part}%
  }{}{\PackageError{exam-gradescope}{Could not patch exam's \string\process@part}{}}%
}{}{\PackageError{exam-gradescope}{Could not patch the parts environment}{}}
%</package>
%    \end{macrocode}
%
% |\@readpoints|[pts] runs for questions, parts and subparts. Points may be macros, or use exam's |\half|.
%
%    \begin{macrocode}
%<*package>
\pretocmd\@readpoints{%
  \begingroup
    \def\half{.5}%
    \protected@edef\eg@points{#1}%
    \directlua{examgradescope.points("\luaescapestring{\@queslevel}",
      "\luaescapestring{\expandafter\detokenize\expandafter{\eg@points}}")}%
  \endgroup
}{}{\PackageError{exam-gradescope}{Could not patch exam's \string\@readpoints}{}}
%</package>
%    \end{macrocode}
%
% The first |\label| in a question names it.
%
%    \begin{macrocode}
%<*package>
\AddToHookWithArguments{label}{%
  \protected@edef\eg@label{#1}%
  \directlua{examgradescope.label("\luaescapestring{\expandafter\detokenize\expandafter{\eg@label}}")}}
%</package>
%    \end{macrocode}
%
% Choices: after each environment defines |\choice| and |\CorrectChoice|, wrap them to count the choices (checkboxes have no counter) and note the correct ones.
%
%    \begin{macrocode}
%<*package>
\newcommand\eg@wrapchoices[1]{%
  \directlua{examgradescope.choices("#1")}%
  \let\eg@choice\choice
  \let\eg@CorrectChoice\CorrectChoice
  \def\choice{\directlua{examgradescope.choice(false)}\eg@choice}%
  \def\CorrectChoice{\directlua{examgradescope.choice(true)}\eg@CorrectChoice}%
  \let\correctchoice\CorrectChoice
}
%</package>
%    \end{macrocode}
%
% exam's own environments are patched. With its overload option, exam-randomizechoices replaces them by environments (taking an optional argument) that call the originals, saved as |\erc@old<name>|; if it was loaded first, those originals are the ones to patch.
%
%    \begin{macrocode}
%<*package>
\newcommand\eg@patchchoices[2]{%
  \ifcsdef{erc@old#1}%
    {\expandafter\apptocmd\csname erc@old#1\endcsname{\eg@wrapchoices{#2}}{}{\eg@patchfailed{#1}}}%
    {\expandafter\apptocmd\csname #1\endcsname{\eg@wrapchoices{#2}}{}{\eg@patchfailed{#1}}}}
\newcommand\eg@patchfailed[1]{%
  \PackageError{exam-gradescope}{Could not patch the #1 environment}{}}
\eg@patchchoices{choices}{mc}
\eg@patchchoices{oneparchoices}{mc}
\eg@patchchoices{checkboxes}{cata}
\eg@patchchoices{oneparcheckboxes}{cata}

\NewDocumentCommand\gradescopequestion{m}{%
  \directlua{examgradescope.question_options("\luaescapestring{\detokenize{#1}}")}}
%</package>
%    \end{macrocode}
%
% \subsubsection{Name and ID regions}
%
%    \begin{macrocode}
%<*package>
\NewDocumentCommand\GradescopeRegion{m +m}{%
  \leavevmode
  \setbox\z@=\hbox{#2}%
  \expanded{\noexpand\latelua{examgradescope.region("#1", \number\wd\z@, \number\ht\z@, \number\dp\z@)}}%
  \box\z@
}
%</package>
%    \end{macrocode}
%
% \subsubsection{Rubrics}
%
% A rubric's source is read as text: descriptions go to Gradescope as written, with math as |$$...$$| and |\textbf| as |**...**|.
%
%    \begin{macrocode}
%<*package>
\NewDocumentEnvironment{rubric}{O{} +b}{%
  \directlua{examgradescope.rubric("\luaescapestring{\detokenize{#1}}", "\luaescapestring{\detokenize{#2}}")}%
  \eg@showrubric
}{}

\NewDocumentCommand\DeclareRubric{O{} m +m}{%
  \directlua{examgradescope.declare_rubric("\luaescapestring{#2}", "\luaescapestring{\detokenize{#1}}",
    "\luaescapestring{\detokenize{#3}}")}}

\NewDocumentCommand\userubric{m +O{}}{%
  \directlua{examgradescope.use_rubric("\luaescapestring{#1}", "\luaescapestring{\detokenize{#2}}")}%
  \eg@showrubric
}
%</package>
%    \end{macrocode}
%
% showrubric: print the rubric when answers are shown.
%
%    \begin{macrocode}
%<*package>
\newcommand\eg@showrubric{%
  \ifeg@show\ifprintanswers
    \par\directlua{examgradescope.show_rubric()}%
  \fi\fi}
\newenvironment{examgradescoperubric}[1]{%
  \list{}{\small
    \def\makelabel##1{\hss\llap{\sffamily##1}}%
    \setlength\itemsep{0pt}\setlength\parsep{0pt}\setlength\topsep{2pt}}%
  \ifx\relax#1\relax\else\item[]\textsf{Rubric (#1 scoring)}\fi
}{\endlist}
\newcommand\examgradescopegroup[2]{%
  \textsf{#1}\ifx\relax#2\relax\else\ (\textit{#2})\fi}
%</package>
%    \end{macrocode}
%
% \subsubsection{Page layout, files and debugregions}
%
%    \begin{macrocode}
%<*package>
\AtBeginDocument{%
  \directlua{examgradescope.set_layout(
    \number\dimexpr 1in + \hoffset + \oddsidemargin\relax,
    \number\dimexpr 1in + \voffset + \topmargin + \headheight + \headsep\relax,
    \number\textwidth, \number\textheight)}%
}

\AtEndDocument{%
  \directlua{
    examgradescope.options.outline = \ifeg@outline true\else false\fi;
    examgradescope.options.rubric = \ifeg@rubric true\else false\fi;
    examgradescope.options.key = \ifeg@key true\else false\fi;
    examgradescope.options.manifest = \ifeg@manifest true\else false\fi;
  }%
}
\directlua{
  luatexbase.add_to_callback("stop_run", function() examgradescope.write_files(tex.jobname) end,
    "examgradescope")
}

\newcommand\examgradescopebox[3]{%
  \color{red}\setlength\fboxsep{0pt}\setlength\fboxrule{0.8pt}%
  \fbox{\parbox[b][#2][t]{#1}{\hfill\colorbox{red}{\color{white}\tiny\ttfamily #3}}}}
\ifeg@debug
  \RequirePackage{xcolor}
  \directlua{examgradescope.load_previous(tex.jobname)}
  \AddToHook{shipout/foreground}{%
    \directlua{examgradescope.draw_page(status.total_pages + 1, tex.pagewidth, tex.pageheight)}}
\fi
%</package>
%    \end{macrocode}
%
% \subsection{The Lua module}
%
% The package's work is done in Lua: recording questions as \TeX{} reads
% them, positions as pages are shipped out, and the files at the end of the
% run.
%
%    \begin{macrocode}
%<*lua>
-- exam-gradescope.lua: record the questions of an exam-class document and
-- write files for Gradescope: an outline, rubrics, a bubble-sheet answer key
-- and an item manifest.

local eg = {}

eg.questions = {}      -- in document order
eg.current = nil       -- the question being read
eg.part = nil          -- the part being read, if any
eg.marks = {}          -- position marks, filled in at shipout
eg.regions = {}        -- identification regions: name, sid
eg.layout = {}         -- text area, in sp from the top left of the page
eg.presets = {}        -- \DeclareRubric
eg.options = {
  mc = "gradescope",               -- where MC and CATA questions go: "gradescope" or "bubble"
  margin = 18 * 65536,             -- added left and right of the text area
  min_height = 12 * 65536,         -- shorter rectangles are dropped
  default_rubric = "",             -- preset for free-response questions with no rubric
  outline = true, rubric = true, key = true, manifest = true,
}

local function warn(message)
  texio.write_nl("term and log", "Package exam-gradescope Warning: " .. message)
  texio.write_nl("term and log", "")
end

local function trim(s)
  return (s:gsub("^%s+", ""):gsub("%s+$", ""))
end

-- ---------------------------------------------------------------------------
-- Recording, while TeX reads the document
-- ---------------------------------------------------------------------------

function eg.start_question(number)
  local q = {number = number, points = 0, parts = {}, correct = {}, choices = 0}
  eg.questions[#eg.questions + 1] = q
  eg.current = q
  eg.part = nil
  return #eg.questions
end

function eg.start_part(number)
  local q = eg.current
  if not q then return 0 end
  local p = {number = number, points = 0}
  q.parts[#q.parts + 1] = p
  eg.part = p
  return #q.parts
end

-- The question, or the part of it, being read.
local function target()
  return eg.part or eg.current
end

-- Points given as \question[pts] or \part[pts]; subparts add to their part.
function eg.points(level, value)
  local n = tonumber(value)
  if not eg.current then return end
  if not n then
    warn("could not read points '" .. value .. "' of question " .. eg.current.number)
    return
  end
  if level == "question" or not eg.part then
    eg.current.points = eg.current.points + n
  else
    eg.part.points = eg.part.points + n
  end
end

-- The first \label in a question names it. Labels with an @ are internal:
-- exam's own (question@3, firstpoints@onpage@1) and exam-randomizechoices'.
function eg.label(name)
  local q = eg.current
  if not q or q.id ~= nil then return end
  if name:find("@", 1, true) then return end
  q.id = name
end

-- kind: "mc" or "cata", from the choices environment.
function eg.choices(kind)
  local q = eg.current
  if q then q.kind = q.kind or kind end
end

-- Choices are lettered A, B, ... in the order they are typeset, which is
-- after any shuffling.
function eg.choice(correct)
  local q = eg.current
  if not q then return end
  q.choices = q.choices + 1
  if correct then q.correct[#q.correct + 1] = string.char(64 + q.choices) end
end

-- \gradescopequestion{type=..., destination=..., title=..., id=..., points=...}
-- In a part, title and points apply to the part.
function eg.question_options(text)
  local q = eg.current
  if not q then return end
  for key, value in pairs(eg.keyvals(text)) do
    if eg.part and (key == "title" or key == "points") then
      if key == "points" then
        eg.part.points_override = tonumber(value)
        if not eg.part.points_override then warn("points should be a number, not '" .. value .. "'") end
      else
        eg.part.title = value
      end
      goto continue
    end
    if key == "type" then
      if value ~= "mc" and value ~= "cata" and value ~= "frq" then
        warn("type should be mc, cata or frq, not '" .. value .. "'")
      else
        q.kind_override = value
      end
    elseif key == "destination" then
      if value ~= "bubble" and value ~= "gradescope" and value ~= "none" then
        warn("destination should be bubble, gradescope or none, not '" .. value .. "'")
      else
        q.destination = value
      end
    elseif key == "title" or key == "id" then
      q[key] = value
    elseif key == "points" then
      q.points_override = tonumber(value)
      if not q.points_override then warn("points should be a number, not '" .. value .. "'") end
    else
      warn("unknown key '" .. key .. "' in \\gradescopequestion")
    end
    ::continue::
  end
end

-- ---------------------------------------------------------------------------
-- Positions, at shipout
-- ---------------------------------------------------------------------------

local function page_number()
  return status.total_pages + 1
end

-- From \latelua: kind is "question", "part" or "end".
function eg.mark(kind, index, part)
  local h, v = pdf.getpos()
  eg.marks[#eg.marks + 1] = {
    kind = kind, index = index, part = part,
    page = page_number(), h = h, y = tex.pageheight - v,
    width = tex.pagewidth, height = tex.pageheight,
  }
end

-- From \latelua before an \hbox of the given size.
function eg.region(name, wd, ht, dp)
  local h, v = pdf.getpos()
  eg.regions[name] = {
    page = page_number(), x1 = h, x2 = h + wd,
    y1 = tex.pageheight - v - ht, y2 = tex.pageheight - v + dp,
    width = tex.pagewidth, height = tex.pageheight,
  }
end

function eg.set_layout(left, top, width, height)
  eg.layout = {left = left, top = top, width = width, height = height}
end

-- ---------------------------------------------------------------------------
-- Small parsers for the LaTeX interface
-- ---------------------------------------------------------------------------

-- The balanced {...} or [...] group starting at position i (which must hold
-- the opening delimiter). Returns its contents and the position after it.
local function group(s, i, open, close)
  open, close = open or "{", close or "}"
  if s:sub(i, i) ~= open then return nil, i end
  local depth = 0
  for j = i, #s do
    local c = s:sub(j, j)
    if c == "\\" then
      -- skip an escaped character
    elseif c == open and s:sub(j - 1, j - 1) ~= "\\" then
      depth = depth + 1
    elseif c == close and s:sub(j - 1, j - 1) ~= "\\" then
      depth = depth - 1
      if depth == 0 then return s:sub(i + 1, j - 1), j + 1 end
    end
  end
  return nil, i
end

local function skip_space(s, i)
  return s:find("[^%s]", i) or #s + 1
end

-- "a=1, b, c = {x, y}" -> {a = "1", b = "true", c = "x, y"}
function eg.keyvals(text)
  local out = {}
  local i = 1
  while i <= #text do
    i = skip_space(text, i)
    local key_end = text:find("[=,]", i) or #text + 1
    local key = trim(text:sub(i, key_end - 1))
    local value = "true"
    i = key_end
    if text:sub(i, i) == "=" then
      i = skip_space(text, i + 1)
      if text:sub(i, i) == "{" then
        value, i = group(text, i)
        i = text:find(",", i) or #text + 1
      else
        local value_end = text:find(",", i) or #text + 1
        value = trim(text:sub(i, value_end - 1))
        i = value_end
      end
    end
    if key ~= "" then out[key] = value end
    i = i + 1
  end
  return out
end

local function weight_of(text, where)
  local a, b = text:match("^%s*([%d.]+)%s*/%s*([%d.]+)%s*$")
  local w = a and tonumber(a) / tonumber(b) or tonumber(trim(text))
  if not w then warn("could not read the weight '" .. text .. "'" .. (where or "")) end
  return w or 0
end

-- Detokenized LaTeX, tidied: one space between words, and no space after a
-- control word (which \detokenize adds) unless a letter follows.
function eg.source(text)
  local s = trim(text:gsub("%s+", " "))
  return (s:gsub("(\\%a+) ([^%a])", "%1%2"))
end

-- LaTeX source of a rubric description, as Gradescope's Markdown: math in
-- $$...$$, \textbf in **...**, \emph and \textit in *...*.
function eg.markdown(text)
  local s = eg.source(text)
  s = s:gsub("\\%( ?", "$$"):gsub(" ?\\%)", "$$")
  s = s:gsub("\\%[ ?", "$$"):gsub(" ?\\%]", "$$")
  -- single $...$ (but not $$...$$) becomes $$...$$
  s = s:gsub("%$%$", "\1"):gsub("%$([^%$\1]-)%$", "\1%1\1"):gsub("\1", "$$")
  local function wrap(command, mark)
    s = s:gsub("\\" .. command .. " ?(%b{})", function(g) return mark .. g:sub(2, -2) .. mark end)
  end
  wrap("textbf", "**")
  wrap("emph", "*")
  wrap("textit", "*")
  s = s:gsub("\\%%", "%%"):gsub("\\&", "&"):gsub("\\#", "#"):gsub("\\_", "_")
  s = s:gsub("~", " "):gsub("\\ldots ?", "…"):gsub("\\dots ?", "…")
  s = s:gsub("``", "“"):gsub("''", "”")
  return trim(s)
end

-- The body of a rubric environment or preset: \item[weight] text, and
-- \begin{rubricgroup}[select one]{description} ... \end{rubricgroup}.
-- Returns a list of entries {item = {description, weight}} or
-- {group = {description, mutually_exclusive, items}}.
function eg.parse_rubric_body(body)
  local entries = {}
  local list = entries        -- where items go: the rubric, or a group's items
  local group_entry = nil
  local i = 1
  local pending = nil         -- the item whose text is being read
  local function finish(upto)
    if pending then
      pending.source = eg.source(body:sub(pending.from, upto - 1))
      pending.description = eg.markdown(pending.source)
      pending.from = nil
      if pending.description == "" then warn("a rubric item has no description") end
      pending = nil
    end
  end
  while true do
    local a = body:find("\\item", i, true)
    local b = body:find("\\begin ?{rubricgroup}", i)
    local c = body:find("\\end ?{rubricgroup}", i)
    local nxt = math.min(a or math.huge, b or math.huge, c or math.huge)
    if nxt == math.huge then
      finish(#body + 1)
      break
    end
    finish(nxt)
    if nxt == a then
      local j = skip_space(body, a + 5)
      local w = "0"
      if body:sub(j, j) == "[" then
        w, j = group(body, j, "[", "]")
      else
        warn("a rubric item has no [weight]; using 0")
      end
      pending = {weight = weight_of(w or "0", " in a rubric item"), from = j}
      list[#list + 1] = pending
    elseif nxt == b then
      local j = body:find("}", b) + 1
      j = skip_space(body, j)
      local options = ""
      if body:sub(j, j) == "[" then options, j = group(body, j, "[", "]") end
      j = skip_space(body, j)
      local description
      description, j = group(body, j)
      local opts = eg.keyvals(options or "")
      group_entry = {
        source = eg.source(description or ""),
        description = eg.markdown(description or ""),
        mutually_exclusive = opts["select one"] ~= nil,
        items = {},
      }
      entries[#entries + 1] = {group = group_entry}
      list = group_entry.items
      i = j
    else
      list = entries
      group_entry = nil
      i = body:find("}", c) + 1
    end
    if nxt == a then i = pending.from end
  end
  -- Items of the top level were collected as bare tables; wrap them.
  for k, entry in ipairs(entries) do
    if not entry.group then entries[k] = {item = entry} end
  end
  return entries
end

-- A preset's entries, with weights scaled and rounded to hundredths.
local function scaled(entries, factor)
  local out = {}
  local function weight(w) return math.floor(w * factor * 100 + 0.5) / 100 end
  local function item(it) return {description = it.description, source = it.source, weight = weight(it.weight)} end
  for _, e in ipairs(entries) do
    if e.item then
      out[#out + 1] = {item = item(e.item)}
    else
      local items = {}
      for _, it in ipairs(e.group.items) do items[#items + 1] = item(it) end
      out[#out + 1] = {group = {description = e.group.description, source = e.group.source,
                                mutually_exclusive = e.group.mutually_exclusive, items = items}}
    end
  end
  return out
end

local function settings_of(text)
  local opts = eg.keyvals(text or "")
  local settings = {scoring_type = "positive", floor = true, ceiling = true}
  if opts.scoring then settings.scoring_type = opts.scoring end
  if opts.floor then settings.floor = opts.floor ~= "false" end
  if opts.ceiling then settings.ceiling = opts.ceiling ~= "false" end
  return settings
end

-- \begin{rubric}[options] body \end{rubric}
function eg.rubric(options, body)
  local t = target()
  if not t then
    warn("a rubric outside any question is ignored")
    return
  end
  t.rubric = {settings = settings_of(options), entries = eg.parse_rubric_body(body)}
end

-- \DeclareRubric[options]{name}{body}; weights are fractions of the points.
function eg.declare_rubric(name, options, body)
  eg.presets[name] = {settings = settings_of(options), entries = eg.parse_rubric_body(body)}
end

-- Find an entry by the start of its description (ignoring case).
local function find_entry(entries, text)
  local want = eg.markdown(text):lower()
  for k, e in ipairs(entries) do
    local d = (e.item or e.group).description:lower()
    if d:sub(1, #want) == want then return k end
  end
  warn("no rubric item starts with '" .. text .. "'")
end

-- \reword{match}{new}, \add[before=match|after=match]{weight}{text}, \remove{match}
local function apply_edits(entries, edits)
  local i = 1
  while true do
    local s, e, command = edits:find("\\(%a+)", i)
    if not s then break end
    local j = skip_space(edits, e + 1)
    if command == "reword" then
      local match, new
      match, j = group(edits, j); j = skip_space(edits, j)
      new, j = group(edits, j)
      local k = match and find_entry(entries, match)
      if k and new then
        local e = entries[k].item or entries[k].group
        e.description, e.source = eg.markdown(new), eg.source(new)
      end
    elseif command == "remove" then
      local match
      match, j = group(edits, j)
      local k = match and find_entry(entries, match)
      if k then table.remove(entries, k) end
    elseif command == "add" then
      local options = ""
      if edits:sub(j, j) == "[" then options, j = group(edits, j, "[", "]"); j = skip_space(edits, j) end
      local w, text
      w, j = group(edits, j); j = skip_space(edits, j)
      text, j = group(edits, j)
      local opts = eg.keyvals(options or "")
      local entry = {item = {description = eg.markdown(text or ""), source = eg.source(text or ""),
                             weight = weight_of(w or "0", " in \\add")}}
      local k = #entries + 1
      if opts.before then
        k = find_entry(entries, opts.before) or k
      elseif opts.after then
        k = (find_entry(entries, opts.after) or #entries) + 1
      end
      table.insert(entries, k, entry)
    else
      warn("unknown rubric edit \\" .. command)
    end
    i = j
  end
end

-- \userubric{name}[edits]: the preset scaled to the points of the question
-- (or part) so far. Returns the rubric, for showing.
function eg.use_rubric(name, edits)
  local t = target()
  local preset = eg.presets[name]
  if not t then return end
  if not preset then
    warn("no rubric named '" .. name .. "'")
    return
  end
  local entries = scaled(preset.entries, t.points_override or t.points)
  apply_edits(entries, edits or "")
  t.rubric = {settings = preset.settings, entries = entries}
end

-- ---------------------------------------------------------------------------
-- Showing a rubric (showrubric)
-- ---------------------------------------------------------------------------

local function number_text(x)
  if x == math.floor(x) then return string.format("%d", x) end
  return (string.format("%.2f", x):gsub("0+$", ""))
end

-- The current question's (or part's) rubric, typeset as a list.
function eg.show_rubric()
  local t = target()
  if not t or not t.rubric then return end
  local function item(it)
    tex.sprint("\\item[" .. number_text(it.weight) .. "]")
    tex.sprint(it.source or it.description)
  end
  tex.sprint("\\begin{examgradescoperubric}{" .. t.rubric.settings.scoring_type .. "}")
  for _, e in ipairs(t.rubric.entries) do
    if e.item then
      item(e.item)
    else
      tex.sprint("\\item[]\\examgradescopegroup{")
      tex.sprint(e.group.source or e.group.description)
      tex.sprint("}{" .. (e.group.mutually_exclusive and "select one" or "") .. "}")
      tex.sprint("\\begin{examgradescoperubric}{}")
      for _, it in ipairs(e.group.items) do item(it) end
      tex.sprint("\\end{examgradescoperubric}")
    end
  end
  tex.sprint("\\end{examgradescoperubric}")
end

-- ---------------------------------------------------------------------------
-- Regions
-- ---------------------------------------------------------------------------

local function round(x)
  return math.floor(x * 10 + 0.5) / 10
end

-- A rectangle on one page, from y1 to y2 (sp from the top), in percent.
local function rect(page, y1, y2, width, height)
  local l = eg.layout
  local x1 = math.max(0, l.left - eg.options.margin)
  local x2 = math.min(width, l.left + l.width + eg.options.margin)
  return {page_number = page, x1 = round(100 * x1 / width), x2 = round(100 * x2 / width),
          y1 = round(100 * y1 / height), y2 = round(100 * y2 / height)}
end

-- The rectangles from one mark to the next, one per page. Slivers shorter
-- than min_height (a mark that fell at the foot of a page) are dropped.
local function span(from, to)
  local l = eg.layout
  local top, bottom = l.top, l.top + l.height
  local rects = {}
  local function add(page, y1, y2)
    if y2 - y1 >= eg.options.min_height then rects[#rects + 1] = rect(page, y1, y2, from.width, from.height) end
  end
  if from.page == to.page then
    add(from.page, from.y, to.y)
  else
    add(from.page, from.y, bottom)
    for page = from.page + 1, to.page - 1 do add(page, top, bottom) end
    add(to.page, top, to.y)
  end
  return rects
end

-- Each question's (and part's) rectangles run from its mark to the next
-- question's (or part's) mark. Call it once every page has been shipped out.
function eg.resolve_regions()
  local starts = {}   -- question and end marks, in document order
  for _, m in ipairs(eg.marks) do
    if m.kind == "question" then
      eg.questions[m.index].mark = m
      starts[#starts + 1] = m
    elseif m.kind == "part" then
      eg.questions[m.index].parts[m.part].mark = m
    else
      starts[#starts + 1] = m
    end
  end
  for i, m in ipairs(starts) do
    local nxt = starts[i + 1]
    -- The last question's answer space is the rest of its page, after its
    -- text, so it runs to the foot of the text area: of the page with the
    -- end of the questions or, if that mark was never shipped out (it was
    -- alone on a page), of the page of the question's last mark.
    if m.kind == "question" and (not nxt or nxt.kind == "end") then
      local last = m
      for _, p in ipairs(eg.questions[m.index].parts) do
        if p.mark and p.mark.page >= last.page then last = p.mark end
      end
      local page = nxt and math.max(nxt.page, last.page) or last.page
      nxt = {page = page, y = eg.layout.top + eg.layout.height, width = m.width, height = m.height}
    end
    if m.kind == "question" and nxt then
      local q = eg.questions[m.index]
      q.crop_rect_list = span(m, nxt)
      for j, p in ipairs(q.parts) do
        if p.mark then
          local after = q.parts[j + 1] and q.parts[j + 1].mark or nxt
          p.crop_rect_list = span(p.mark, after)
        end
      end
    end
  end
end

-- ---------------------------------------------------------------------------
-- Questions as Gradescope sees them
-- ---------------------------------------------------------------------------

local function kind_of(q)
  return q.kind_override or q.kind or "frq"
end

local function destination_of(q)
  if q.destination then return q.destination end
  if kind_of(q) == "frq" then return "gradescope" end
  return eg.options.mc
end

local function questions_to(destination)
  local out = {}
  for _, q in ipairs(eg.questions) do
    if destination_of(q) == destination then out[#out + 1] = q end
  end
  return out
end

-- Points of a part: given by \part[pts] or \gradescopequestion{points=}.
-- If none of a question's parts has any, the question's points are split
-- evenly among them.
local function part_points(q, p)
  if p.points_override then return p.points_override end
  for _, other in ipairs(q.parts) do
    if other.points > 0 or other.points_override then return p.points end
  end
  if not q.warned_split and q.points > 0 then
    warn("the parts of question " .. q.number .. " have no points; splitting its "
      .. q.points .. " evenly (give each part's points with \\gradescopequestion{points=...})")
    q.warned_split = true
  end
  return q.points / #q.parts
end

-- A question with parts is worth the sum of its parts.
local function total_points(q)
  if q.points_override then return q.points_override end
  if #q.parts == 0 then return q.points end
  local sum = 0
  for _, p in ipairs(q.parts) do sum = sum + part_points(q, p) end
  return sum
end

local function part_title(p, k)
  return p.title or ("(" .. string.char(96 + k) .. ")")
end

-- When some questions go to the bubble sheet, Gradescope's numbering differs
-- from the printed one, so the title gives the printed number.
local function title_of(q)
  if q.title then return q.title end
  if #questions_to("bubble") > 0 then return "Question " .. q.number end
  return ""
end

-- The rubric for MC and CATA questions, following the usual conventions:
-- MC: "Correct" for full points and "The correct response is **X**" for 0.
-- CATA: negative scoring, one item per choice deducting points/choices.
local function choice_rubric(q)
  local points = total_points(q)
  local correct = {}
  for _, letter in ipairs(q.correct) do correct[letter] = true end
  local settings = {scoring_type = "positive", floor = true, ceiling = true}
  if kind_of(q) == "mc" then
    return {settings = settings, entries = {
      {item = {description = "Correct", weight = points}},
      {item = {description = "The correct response is **" .. table.concat(q.correct, "") .. "**", weight = 0}},
    }}
  end
  settings.scoring_type = "negative"
  local entries = {}
  local deduction = q.choices > 0 and points / q.choices or 0
  for k = 1, q.choices do
    local letter = string.char(64 + k)
    entries[#entries + 1] = {item = {
      description = letter .. (correct[letter] and " **should** be selected" or " **should not** be selected"),
      weight = deduction,
    }}
  end
  entries[#entries + 1] = {item = {description = "Correct", weight = 0}}
  return {settings = settings, entries = entries}
end

local function rubric_of(q, t)
  if t.rubric then return t.rubric end
  if t == q and kind_of(q) ~= "frq" then return choice_rubric(q) end
  local preset = eg.presets[eg.options.default_rubric]
  local points = t == q and total_points(q) or part_points(q, t)
  if preset then return {settings = preset.settings, entries = scaled(preset.entries, points)} end
end

-- ---------------------------------------------------------------------------
-- JSON
-- ---------------------------------------------------------------------------

-- Objects keep their keys in the order given: obj{"a", 1, "b", 2}.
local object_mt = {}
local function obj(list)
  return setmetatable({pairs = list}, object_mt)
end

local function number_json(x)
  if x == math.floor(x) and math.abs(x) < 1e15 then return string.format("%d", x) end
  return (string.format("%.6f", x):gsub("0+$", ""))
end

local escapes = {['"'] = '\\"', ['\\'] = '\\\\', ['\n'] = '\\n', ['\t'] = '\\t', ['\r'] = '\\r'}
local function string_json(s)
  return '"' .. s:gsub('[%c"\\]', function(c)
    return escapes[c] or string.format("\\u%04x", c:byte())
  end) .. '"'
end

local function to_json(value, indent)
  indent = indent or ""
  local inner = indent .. "  "
  local t = type(value)
  if t == "string" then return string_json(value) end
  if t == "number" then return number_json(value) end
  if t == "boolean" then return tostring(value) end
  if value == nil then return "null" end
  local items = {}
  if getmetatable(value) == object_mt then
    for i = 1, #value.pairs, 2 do
      local v = value.pairs[i + 1]
      if v ~= nil then items[#items + 1] = inner .. string_json(value.pairs[i]) .. ": " .. to_json(v, inner) end
    end
    if #items == 0 then return "{}" end
    return "{\n" .. table.concat(items, ",\n") .. "\n" .. indent .. "}"
  end
  for _, v in ipairs(value) do items[#items + 1] = inner .. to_json(v, inner) end
  if #items == 0 then return "[]" end
  return "[\n" .. table.concat(items, ",\n") .. "\n" .. indent .. "]"
end
eg.to_json = to_json

local function rect_json(r)
  return obj{"page_number", r.page_number, "x1", r.x1, "x2", r.x2, "y1", r.y1, "y2", r.y2}
end

local function rects_json(list)
  local out = {}
  for _, r in ipairs(list or {}) do out[#out + 1] = rect_json(r) end
  return out
end

local function region_json(name)
  local r = eg.regions[name]
  if not r then return nil end
  local function pct(x, total) return round(100 * x / total) end
  return obj{"page_number", r.page, "x1", pct(r.x1, r.width), "x2", pct(r.x2, r.width),
             "y1", pct(r.y1, r.height), "y2", pct(r.y2, r.height)}
end

-- ---------------------------------------------------------------------------
-- Output files
-- ---------------------------------------------------------------------------

function eg.outline()
  local questions = {}
  for _, q in ipairs(questions_to("gradescope")) do
    if #q.parts > 0 then
      local children = {}
      for _, p in ipairs(q.parts) do
        children[#children + 1] = obj{"title", part_title(p, #children + 1), "weight", part_points(q, p),
                                      "crop_rect_list", rects_json(p.crop_rect_list)}
      end
      questions[#questions + 1] = obj{"title", title_of(q), "children", children}
    else
      questions[#questions + 1] = obj{"title", title_of(q), "weight", total_points(q),
                                      "crop_rect_list", rects_json(q.crop_rect_list)}
    end
  end
  return obj{"identification_regions", obj{"name", region_json("name"), "sid", region_json("sid")},
             "question_data", questions}
end

-- One rubric in the gradescope-tampermonkey importer's format. Positions are
-- given only when there are both ungrouped items and groups.
local function rubric_json(number, title, rubric)
  local has_items, has_groups = false, false
  for _, e in ipairs(rubric.entries) do
    if e.item then has_items = true else has_groups = true end
  end
  local mixed = has_items and has_groups
  local items, groups = {}, {}
  for position, e in ipairs(rubric.entries) do
    local pos = mixed and position - 1 or nil
    if e.item then
      items[#items + 1] = obj{"description", e.item.description, "weight", e.item.weight, "position", pos}
    else
      local members = {}
      for _, it in ipairs(e.group.items) do
        members[#members + 1] = obj{"description", it.description, "weight", it.weight}
      end
      groups[#groups + 1] = obj{"description", e.group.description,
                                "mutually_exclusive", e.group.mutually_exclusive, "position", pos,
                                "rubric_items", members}
    end
  end
  local s = rubric.settings
  return obj{"question", number, "title", title ~= "" and title or nil,
             "scoring_type", s.scoring_type, "floor", s.floor, "ceiling", s.ceiling,
             "rubric_items", items, "rubric_item_groups", #groups > 0 and groups or nil}
end

function eg.rubrics()
  local out = {}
  for n, q in ipairs(questions_to("gradescope")) do
    if #q.parts > 0 then
      for k, p in ipairs(q.parts) do
        local r = rubric_of(q, p)
        if r then out[#out + 1] = rubric_json(n .. "." .. k, part_title(p, k), r) end
      end
    else
      local r = rubric_of(q, q)
      if r then out[#out + 1] = rubric_json(tostring(n), title_of(q), r) end
    end
  end
  return obj{"questions", out}
end

-- Like aux_to_csv.lua, points always have a decimal point.
local function csv_points(x)
  local s = number_json(x)
  if not s:find("%.") then s = s .. ".0" end
  return s
end

function eg.key_csv()
  local rows = {"Question,Answers,Points,Type"}
  for _, q in ipairs(questions_to("bubble")) do
    rows[#rows + 1] = string.format("%d,%s,%s,Exact", q.number, table.concat(q.correct, ""), csv_points(total_points(q)))
  end
  return table.concat(rows, "\n") .. "\n"
end

function eg.manifest()
  local items = {}
  for _, q in ipairs(eg.questions) do
    local parts = nil
    if #q.parts > 0 then
      parts = {}
      for _, p in ipairs(q.parts) do parts[#parts + 1] = obj{"points", part_points(q, p)} end
    end
    items[#items + 1] = obj{"number", q.number, "id", q.id or tostring(q.number), "points", total_points(q),
                            "type", kind_of(q), "destination", destination_of(q),
                            "answer", #q.correct > 0 and table.concat(q.correct, "") or nil,
                            "parts", parts}
  end
  return obj{"items", items}
end

local function write_file(path, text)
  local f, err = io.open(path, "w")
  if not f then
    warn("could not write " .. path .. ": " .. tostring(err))
    return
  end
  f:write(text)
  f:close()
  texio.write_nl("term and log", "exam-gradescope: wrote " .. path)
end

-- ---------------------------------------------------------------------------
-- debugregions: the previous run's rectangles, drawn on each page
-- ---------------------------------------------------------------------------

local function write_debug(jobname)
  local lines = {"return {"}
  local function add(list, label)
    for _, r in ipairs(list or {}) do
      lines[#lines + 1] = string.format("  {%d, %s, %s, %s, %s, %q},", r.page_number, r.x1, r.x2, r.y1, r.y2, label)
    end
  end
  for n, q in ipairs(questions_to("gradescope")) do
    if #q.parts > 0 then
      for k, p in ipairs(q.parts) do add(p.crop_rect_list, n .. "." .. k) end
    else
      add(q.crop_rect_list, tostring(n))
    end
  end
  for _, name in ipairs({"name", "sid"}) do
    local r = region_json(name)
    if r then
      local v = {}
      for i = 1, #r.pairs, 2 do v[r.pairs[i]] = r.pairs[i + 1] end
      add({v}, name)
    end
  end
  lines[#lines + 1] = "}"
  local f = io.open(jobname .. ".egregions", "w")
  if f then
    f:write(table.concat(lines, "\n"), "\n")
    f:close()
  end
end

eg.previous = nil
function eg.load_previous(jobname)
  local chunk = loadfile(jobname .. ".egregions")
  eg.previous = chunk and chunk() or {}
end

-- picture-mode commands drawing the previous run's rectangles on a page.
-- In the shipout/foreground hook the origin is the top left of the page and
-- the unit is 1pt, so a box sits at (x1, -y2) and rises to y1.
function eg.draw_page(page, width, height)
  if not eg.previous then return end
  for _, r in ipairs(eg.previous) do
    if r[1] == page then
      local x1, x2 = r[2] * width / 100, r[3] * width / 100
      local y1, y2 = r[4] * height / 100, r[5] * height / 100
      tex.sprint(string.format(
        "\\put(%.2f,%.2f){\\examgradescopebox{%.2fpt}{%.2fpt}{%s}}",
        x1 / 65536, -y2 / 65536, (x2 - x1) / 65536, (y2 - y1) / 65536, r[6]))
    end
  end
end

function eg.write_files(jobname)
  eg.resolve_regions()
  write_debug(jobname)
  local o = eg.options
  if o.key and #questions_to("bubble") > 0 then write_file(jobname .. ".key.csv", eg.key_csv()) end
  if #questions_to("gradescope") > 0 then
    if o.outline then write_file(jobname .. ".outline.json", to_json(eg.outline()) .. "\n") end
    local rubrics = eg.rubrics()
    if o.rubric and #rubrics.pairs[2] > 0 then write_file(jobname .. ".rubric.json", to_json(rubrics) .. "\n") end
  end
  if o.manifest and #eg.questions > 0 then write_file(jobname .. ".json", to_json(eg.manifest()) .. "\n") end
end

return eg
%</lua>
%    \end{macrocode}
%
% \Finale
\endinput
